\scanpage{51}
\mspage{4}
% archive foliation: 50
Soit $\Delta$ un diviseur étale de degré $6$ dans $X/S$, muni d'une structure octaédrale i.e. d'une \deleted{\unsure{flèche}} \deleted{de degré $2$} \inserted{opposition involutive}
\[
\underline a_\Delta:\Delta\longrightarrow\Delta
\qquad\deleted{\Delta\longrightarrow\Delta'}
\]
\deleted{où $\Delta'$ est fini étale de degré $3$ sur $S$.}
Correspondant à une opération libre de $\{1,a_\Delta\}$ sur $\Delta/S$. On dit que cette structure octaédrale est \deleted{\uline{homogène}, \unsure{dans}} \inserted{\uline{régulière} dans} $X$, si le groupe
\[
\Gamma=\underline{\operatorname{Aut}}^+_{\mathrm{oct}}(\Delta)
\]
opère sur $X$, de façon compatible avec les op. sur $\Delta$ (cette op. de $\Gamma$ sur $X$ est alors unique).

\margin{Il faut un petit argument pour \unsure{montrer} que tout autom. de la structure oct. $\Delta$ induit par un autom. de $(X,\Delta)$ est direct : OPS que l'autom. fixe les points de $\{0,\infty\}=\Delta_1$, et on se trouve dans les autom. de \unsure{``ceinture''} par $\mu_4$, qui sont bien directs\ldots}

\uline{Pour} $\Delta$ \inserted{diviseur} fixé, on voit qu'il y a au plus une structure octaédrale sur $\Delta$ qui soit \uline{régulière}, et il y en a une ssi $\Delta$ est une ``configuration octaédrale''. Si $\Delta$ est une telle configuration, on vient de voir \inserted{déjà} qu'elle est munie d'une \deleted{unique} structure octaédrale régulière (et \uline{tout} autom. de $X$ qui \deleted{fixe} \inserted{conserve} $\Delta$ \deleted{fixe} \inserted{conserve} cette structure octaédrale). Inversement, si $\Delta$ est muni d'une str. oct. rég, alors le sous-groupe $\underline V$ de $\underline\Gamma=\underline{\operatorname{Aut}}_{\mathrm{oct}}\deleted{\unsure{\mathrm{oc}}}(\Delta)$ inv. abélien non trivial est un groupe loc. isom. à $\mathbb D_{2S}$, qui opère fidèlement sur $X$, et on voit aussitôt que le

\margin{Si on a une configuration octaédrale régulière \deleted{\unsure{$\Delta$}} dans un \deleted{\unsure{fibré}} \inserted{\unsure{en} droites projectives} $X$, alors le revêtement des orientations de \deleted{\unsure{$\Delta$}} est can. isom. au revêtement $\mu_4^*\unsure{=}\operatorname{Spec}\mathcal O_S[i]/(i^2+1)$ de $S$\ldots}
